% ============================================================
%  电商客服 Agentic RAG — 面试演示
%  编译：Overleaf / XeLaTeX
% ============================================================
\documentclass[aspectratio=169,10pt]{beamer}

%% ── 中文支持 ─────────────────────────────────────────────
\usepackage[UTF8]{ctex}   % Overleaf / XeLaTeX

%% ── 工具包 ───────────────────────────────────────────────
\usepackage{tcolorbox}
\tcbuselibrary{skins,breakable}
\usepackage{booktabs}
\usepackage{colortbl}
\usepackage{array}
\usepackage{tabularx}
\usepackage{tikz}
\usetikzlibrary{arrows.meta,positioning,shapes.geometric}
\usepackage{multirow}
\usepackage{amsmath}

%% ── 颜色 ─────────────────────────────────────────────────
\definecolor{Cblue}{RGB}{59,110,182}
\definecolor{Cbluelt}{RGB}{232,240,252}
\definecolor{Cgreelt}{RGB}{228,246,234}
\definecolor{Camberlt}{RGB}{254,243,218}
\definecolor{Cgraylt}{RGB}{245,246,250}
\definecolor{Cpurplelt}{RGB}{242,237,252}
\definecolor{Ctextdark}{RGB}{20,30,55}
\definecolor{Ctextmain}{RGB}{45,55,75}
\definecolor{Cmuted}{RGB}{120,130,150}
\definecolor{Cred}{RGB}{185,40,40}
\definecolor{Cgreen}{RGB}{30,140,60}
\definecolor{Ctakebg}{RGB}{20,30,55}

%% ── Beamer 设置 ──────────────────────────────────────────
\usetheme{default}
\usefonttheme{professionalfonts}
\setbeamertemplate{navigation symbols}{}
\setbeamercolor{normal text}{fg=Ctextmain,bg=white}
\setbeamercolor{background canvas}{bg=white}
\setbeamercolor{frametitle}{fg=Ctextdark,bg=white}

%% 帧标题：节标签 + 大标题 + 副标题
\setbeamertemplate{frametitle}{%
  \vspace*{0.18cm}%
  \noindent\begin{beamercolorbox}[wd=\dimexpr\paperwidth-1.2cm\relax,leftskip=0.8cm,rightskip=0.4cm]{frametitle}%
    {\fontsize{6.5}{8}\selectfont\bfseries\color{Cblue}%
      \MakeUppercase{\insertsectionhead}}\\[0.06cm]%
    {\fontsize{13}{15}\selectfont\bfseries\color{Ctextdark}\insertframetitle}%
    \ifx\insertframesubtitle\@empty\else%
      \\[0.03cm]{\small\color{Cmuted}\insertframesubtitle}%
    \fi%
  \end{beamercolorbox}%
  \vspace{0.06cm}%
  \noindent\makebox[\paperwidth][l]{\hspace*{0.8cm}{\color{Cblue!20}\rule{\dimexpr\paperwidth-1.2cm\relax}{0.5pt}}}%
  \vspace{0.08cm}%
}

%% 页脚：右下角页码
\setbeamertemplate{footline}{%
  \hfill{\tiny\color{Cmuted}\insertframenumber\,/\,\inserttotalframenumber}%
  \hspace{5pt}\vspace{3pt}%
}

%% 列表符号
\setbeamertemplate{itemize item}{\color{Cblue}\small\textbullet}
\setbeamertemplate{itemize subitem}{\color{Cblue!70}\footnotesize$\circ$}

%% ── tcolorbox 预设 ───────────────────────────────────────
\tcbset{
  qbase/.style={
    boxrule=0.4pt, arc=4pt,
    left=6pt, right=6pt, top=4pt, bottom=4pt,
    fonttitle=\small\bfseries
  },
  tcgray/.style  ={qbase, colback=Cgraylt,   colframe=Cgraylt!60!Cmuted},
  tcblue/.style  ={qbase, colback=Cbluelt,   colframe=Cblue!30,   coltitle=Cblue},
  tcgreen/.style ={qbase, colback=Cgreelt,   colframe=Cgreen!30,  coltitle=Cgreen},
  tcamber/.style ={qbase, colback=Camberlt,  colframe=orange!30,  coltitle=orange!70!black},
  tcpurple/.style={qbase, colback=Cpurplelt, colframe=purple!20,  coltitle=purple!60},
  tcred/.style   ={qbase, colback=Cred!6,    colframe=Cred!25,    coltitle=Cred},
}

%% Takeaway 命令（全宽深色条）
\newcommand{\takeaway}[1]{%
  \vfill%
  \begin{tcolorbox}[
    colback=Ctakebg, colframe=Ctakebg,
    arc=0pt, boxrule=0pt,
    left=7pt, right=7pt, top=4pt, bottom=4pt,
    width=\linewidth, nobeforeafter
  ]%
    {\scriptsize\color{white}\textbf{Takeaway}\enspace%
     \color{white!88!Ctakebg}#1}%
  \end{tcolorbox}%
}

%% 辅助命令
\newcommand{\hi}[1]{{\bfseries\color{Cgreen}#1}}      % 好数字 / 强调
\newcommand{\bad}[1]{{\bfseries\color{Cred}#1}}       % 差数字
\newcommand{\mono}[1]{{\ttfamily\small\color{Cblue!90!black}#1}}  % 代码
\newcommand{\cnum}[1]{%
  \tikz[baseline=(n.base)] \node[
    circle,
    draw=Cmuted!70,
    text=Ctextmain,
    line width=0.3pt,
    inner sep=0.35pt,
    minimum size=1.55em,
    font=\scriptsize\bfseries
  ] (n) {#1};%
}

\setlength{\leftmargini}{1em}

%% ============================================================
\begin{document}
%% ============================================================


%% ── P1 个人定位 ─────────────────────────────────────────
\section{开篇·定位}

\begin{frame}[c]
\begin{center}

{\Large\bfseries\color{Ctextdark} 电商客服 Agentic RAG 系统}\\[0.2cm]
{\small\color{Cmuted} NTU $\cdot$ LLM 算法 / RAG 方向}\\[0.4cm]

\begin{tcolorbox}[tcblue, width=0.88\textwidth, center title, halign=center]
  \small
  {\bfseries\color{Cblue}Hybrid Retrieval}\enspace·\enspace
  {\bfseries\color{Cblue}Retrieval Eval}\enspace·\enspace
  {\bfseries\color{Cgreen}Agentic RAG}\enspace·\enspace
  {\bfseries\color{Cgreen}Knowledge Grounding}\enspace·\enspace
  {\bfseries\color{orange!80!black}Ablation Study}
\end{tcolorbox}

\vspace{0.35cm}

\begin{tcolorbox}[tcgray, width=0.88\textwidth, center title]
  \small
  这套工作按一条比较标准的 RAG 工程路径推进：\\[0.2em]
  先定义 gold-doc 指标，再通过\textbf{负结果}和失败案例定位瓶颈，\\
  最后用 \textbf{targeted fix + ablation} 验证改动是否有效。
\end{tcolorbox}

\end{center}
\takeaway{重点是把检索系统做成可测、可量化、可复盘的工程。}
\end{frame}


%% ── P2 项目定位 ─────────────────────────────────────────
\begin{frame}
\frametitle{项目定位}
\framesubtitle{与普通 RAG demo 的本质区别}

\begin{columns}[T]
\column{0.64\textwidth}
{\small
\begin{tabularx}{\linewidth}{lXXX}
\toprule
& {\color{Cmuted}普通 RAG demo} & {\color{Cmuted}传统规则客服} & {\bfseries\color{Cblue}本系统} \\
\midrule
\rowcolor{Cgraylt}
召回 & 无量化评测 & 精确匹配 & \hi{Recall@k / MRR} \\
回答 & 生成即发出 & 模板 & \hi{生成→验证→决策} \\
\rowcolor{Cgraylt}
失败 & 幻觉 & 转人工 & \hi{grounding 闸门} \\
迭代 & 无 & 手工维护 & \hi{SupportCase 飞轮} \\
\bottomrule
\end{tabularx}}

\vspace{0.2cm}
\begin{tcolorbox}[tcblue]
\small 系统输出不是一段回答文本，而是一条可复盘的 \mono{SupportCase}：\\
意图、证据、引用、置信度、验证结果、失败原因和最终动作都会被记录下来。
\end{tcolorbox}

\column{0.33\textwidth}
\begin{tcolorbox}[tcgray, title={\color{Ctextdark}三条设计原则}]
\small
\begin{itemize}\setlength\itemsep{0.3em}
  \item \textbf{LLM} propose → 工具 verify
  \item 失败不浪费 → 飞轮积累
  \item 可测才可信 → 每步有数字
\end{itemize}
\end{tcolorbox}
\end{columns}

\takeaway{定位成可验证、可审计的客服决策系统，而非 RAG demo。}
\end{frame}


%% ── P3 痛点 × 架构回应 ──────────────────────────────────
\begin{frame}
\frametitle{真实痛点 × 架构回应}
\framesubtitle{过渡页——快讲}

{\footnotesize
\begin{tabularx}{\textwidth}{>{\bfseries}lX}
\toprule
客服痛点 & 系统级回应 \\
\midrule
\rowcolor{Cgraylt}
\cnum{1} 幻觉报错价格 / 参数 &
  Grounding + 数字引用检查 + 一致性验证 \\
\cnum{2} 答非所问 &
  意图路由六分类（路由准确率 100\% on 28 eval）\\
\rowcolor{Cgraylt}
\cnum{3} ``600以内''买到 \bad{899元} 产品 &
  预算解析 + post-retrieval 价格过滤
  \quad{\bfseries\color{Cgreen}[NEW]} \\
\cnum{4} ``A和B哪个好''遗漏一半候选 &
  Compound query 分解 + 双实体保底注入
  \quad{\bfseries\color{Cgreen}[NEW]} \\
\bottomrule
\end{tabularx}}

\vspace{0.15cm}
\begin{tcolorbox}[tcblue]
\footnotesize
\cnum{3} 和 \cnum{4} 来自 \textbf{FailureMemory} 对历史失败案例的归纳，不是预先列出的功能项。\\
这里想说明的是：先把问题测出来，再决定修什么。
\end{tcolorbox}

\takeaway{系统设计围绕客服风险展开，每个模块都对应一个可验证的问题。}
\end{frame}


%% ── P4 整体架构 ─────────────────────────────────────────
\section{技术深挖}

\begin{frame}
\frametitle{整体架构}
\framesubtitle{端到端流程}

\begin{columns}[T]
\column{0.50\textwidth}
\vspace{0.1cm}
\begin{tikzpicture}[
  node distance=0.18cm,
  b/.style={draw=Cblue!25, fill=Cgraylt, rounded corners=3pt,
            font=\fontsize{7}{8.5}\selectfont, minimum height=0.44cm,
            minimum width=4.6cm, align=center, text=Ctextmain},
  nb/.style={b, draw=Cgreen!40, fill=Cgreelt,
             text=Cgreen!80!black, font=\fontsize{7}{8.5}\selectfont\bfseries},
  vb/.style={b, draw=orange!40, fill=Camberlt},
  mb/.style={b, draw=purple!25, fill=Cpurplelt},
  db/.style={b, draw=Cred!25,   fill=Cred!6},
  arr/.style={-{Stealth[length=2.5pt]}, Cmuted, line width=0.45pt}
]
\node[b]  (q)   {用户问题};
\node[b,  below=of q]   (ro)  {意图路由（6 类）};
\node[b,  below=of ro]  (ret) {混合检索（dense + BM25 + RRF，父子文档）};
\node[nb, below=0.08cm of ret, xshift=-1.15cm] (pf) {+ 价格约束过滤};
\node[nb, below=0.08cm of ret, xshift=+1.15cm] (cd) {+ 复合查询分解};
\node[b,  below=0.5cm of ret]  (gen) {生成（LLM / catalog 模板）};
\node[vb, below=of gen]  (ver) {三层验证（grounding / 引用 / 一致性）};
\node[db, below=of ver]  (dec) {置信决策（ok / caution / handoff）};
\node[mb, below=of dec]  (sc)  {SupportCase 落库};
\node[mb, below=of sc]   (fm)  {FailureMemory → CaseMemory → 记忆先验};
\draw[arr](q)--(ro); \draw[arr](ro)--(ret);
\draw[arr](ret)--(gen); \draw[arr](gen)--(ver);
\draw[arr](ver)--(dec); \draw[arr](dec)--(sc); \draw[arr](sc)--(fm);
\end{tikzpicture}

\column{0.46\textwidth}
\begin{tcolorbox}[tcgray, title={\color{Ctextdark}设计原则}]
\small\begin{itemize}\setlength\itemsep{0.3em}
  \item \textbf{LLM} 负责 propose
  \item \textbf{确定性工具} 负责 verify
  \item \textbf{记忆} 负责 accumulate
\end{itemize}
\end{tcolorbox}
\vspace{0.25cm}
\begin{tcolorbox}[tcgray, title={\color{Ctextdark}语料规模}]
\small\begin{itemize}\setlength\itemsep{0.2em}
  \item 40 商品 · 5 政策
  \item 205 chunks（NSCC 重建索引）
  \item 28 题 eval set（gold 标注）
\end{itemize}
\end{tcolorbox}
\end{columns}

\takeaway{LLM 负责 propose，确定性工具负责 verify，记忆负责 accumulate。}
\end{frame}


%% ── P5 混合检索 ─────────────────────────────────────────
\begin{frame}
\frametitle{混合检索}
\framesubtitle{为什么 dense alone 不够}

\begin{columns}[T]
\column{0.31\textwidth}
\begin{tcolorbox}[tcblue, title={\color{Cblue}Dense（语义）}]
\small\begin{itemize}\setlength\itemsep{0.25em}
  \item paraphrase-multilingual-\\MiniLM-L12-v2
  \item 余弦相似度闸门 \mono{0.35}
  \item 擅长：口语 / 语义改写
  \item {\small\color{Cmuted}型号数字退化为随机位置}
\end{itemize}
\end{tcolorbox}

\column{0.31\textwidth}
\begin{tcolorbox}[tcamber, title={\color{orange!80!black}BM25（词法）}]
\small\begin{itemize}\setlength\itemsep{0.25em}
  \item jieba 分词
  \item IDF 权重对精确匹配极敏感
  \item 擅长：\mono{"H900"}、\mono{"3000mAh"}
  \item {\small\color{Cmuted}``H900'' 不在训练语料里}
\end{itemize}
\end{tcolorbox}

\column{0.34\textwidth}
\begin{tcolorbox}[tcgreen, title={\color{Cgreen}RRF + 父子文档}]
\small\begin{itemize}\setlength\itemsep{0.25em}
  \item $\tfrac{1}{k+\mathrm{rank}}$ 加权融合
  \item 无需调相似度绝对值 scale
  \item 检索细粒度子 chunk
  \item 喂 LLM 完整父卡（保留上下文）
\end{itemize}
\end{tcolorbox}
\end{columns}

\vspace{0.25cm}

\begin{columns}[c]
\column{0.22\textwidth}
\begin{tcolorbox}[tcgray, halign=center]
{\Large\bfseries\color{Cblue}40}\\{\small 商品}
\end{tcolorbox}
\column{0.22\textwidth}
\begin{tcolorbox}[tcgray, halign=center]
{\Large\bfseries\color{Cblue}205}\\{\small chunks}
\end{tcolorbox}
\column{0.27\textwidth}
\begin{tcolorbox}[tcgray, halign=center]
{\Large\bfseries\color{Cblue}7}\\{\small 款耳机（混淆簇）}
\end{tcolorbox}
\column{0.27\textwidth}
\begin{tcolorbox}[tcgray, halign=center]
{\Large\bfseries\color{Cblue}4}\\{\small 款键盘（混淆簇）}
\end{tcolorbox}
\end{columns}

\takeaway{口语问题靠 dense，型号参数靠 BM25，RRF 兼顾两者。}
\end{frame}


%% ── P6 生成后验证 ───────────────────────────────────────
\begin{frame}
\frametitle{生成后验证}
\framesubtitle{三层可信度闸门}

\begin{columns}[T]
\column{0.32\textwidth}
\begin{tcolorbox}[tcblue, title={\cnum{1} 句级 Grounding}]
\small
\mono{encode(回答句)} $\leftrightarrow$ \mono{encode(context)}\\[0.2em]
余弦相似度 $> 0.42$ 为有支撑\\
支撑率 $< 50\%$
→ {\bfseries\color{orange}caution}
\end{tcolorbox}

\column{0.32\textwidth}
\begin{tcolorbox}[tcamber, title={\cnum{2} 数字引用检查}]
\small
regex 检测价格 / 规格数字\\[0.2em]
验证是否有 \mono{[资料N]} 标注\\
无引用 → {\bfseries\color{orange}caution}
\end{tcolorbox}

\column{0.32\textwidth}
\begin{tcolorbox}[tcpurple, title={\cnum{3} LLM 一致性}]
\small
判断是否「矛盾 / 资料外 / 通过」\\（有 API key 时启用）\\[0.2em]
矛盾 → {\bfseries\color{Cred}handoff}
\end{tcolorbox}
\end{columns}

\vspace{0.25cm}
\begin{tcolorbox}[tcgray]
\centering\small
三层 → \mono{final\_decision(grounding,\,consistency,\,citation)}
→ \hi{ok} \,/\, {\bfseries\color{orange}caution} \,/\, \bad{handoff}
\end{tcolorbox}

\vspace{0.15cm}
\begin{tcolorbox}[tcblue]
\small\textbf{无 LLM key 时}：grounding 和引用检查仍可运行（确定性计算），系统不依赖 LLM 自评。
\end{tcolorbox}

\takeaway{回答返回给用户之前，要先通过可信度校验。}
\end{frame}


%% ── P7 评测体系 ─────────────────────────────────────────
\begin{frame}
\frametitle{评测体系}
\framesubtitle{科学方法论的基石}

\begin{columns}[T]
\column{0.50\textwidth}
\begin{tcolorbox}[tcgray, title={\color{Ctextdark}评测集设计}]
\small\begin{itemize}\setlength\itemsep{0.3em}
  \item \textbf{28 题} gold 标注：\\
    \mono{expected\_intent} / \mono{expected\_action}\\
    / \mono{gold\_doc\_ids}
  \item 刻意构造混淆簇：7 款降噪耳机、\\4 款机械键盘、5 款杯类
  \item 指标为 query \textbf{macro-average}；\\所有 28 题均有 gold，全部参与计算
\end{itemize}
\end{tcolorbox}
\vspace{0.2cm}
\begin{tcolorbox}[tcgray, title={\color{Ctextdark}执行环境}]
\small\begin{itemize}\setlength\itemsep{0.2em}
  \item NSCC PBS CPU 节点，\mono{HF\_HUB\_OFFLINE=1}
  \item env flag 控制条件，5 条件共用单 retriever
\end{itemize}
\end{tcolorbox}

\column{0.46\textwidth}
\begin{tcolorbox}[tcgreen, title={\color{Cgreen}Baseline 结果}]
\vspace{0.1cm}\centering
\begin{tabular}{cccc}
{\large\bfseries\color{Cgreen}0.907} &
{\large\bfseries\color{Cgreen}0.981} &
{\large\bfseries\color{Cgreen}0.981} &
{\large\bfseries\color{Cgreen}0.981}\\
{\tiny R@1} & {\tiny R@3} & {\tiny R@5} & {\tiny MRR}
\end{tabular}
\end{tcolorbox}
\vspace{0.2cm}
\begin{tcolorbox}[tcblue]
\small\textbf{为什么用 Recall@k 而非 RAGAS？}\\[0.2em]
gold-doc 召回是一阶问题；生成质量需要 LLM key，可以加，但先把检索搞对是前提。
\end{tcolorbox}
\vspace{0.15cm}
\begin{tcolorbox}[tcgray]
\small\textbf{recall@1=0.907 怎么算？}\\
28 题，每题看 retrieved\_doc\_ids[0] 是否在 gold\_doc\_ids 里，命中数/28，macro-average。
\end{tcolorbox}
\end{columns}

\takeaway{先有量化指标，才有资格谈「改进了多少」。}
\end{frame}


%% ── P8 Reranker 负结果 ──────────────────────────────────
\section{实验与进化}

\begin{frame}
\frametitle{Honest Evaluation — Reranker 负结果}
\framesubtitle{评测驱动决策}

{\footnotesize
\renewcommand{\arraystretch}{1.18}
\begin{tabularx}{\textwidth}{lcccc>{\raggedright\arraybackslash}X}
\toprule
\bfseries 配置 & \bfseries R@1 & \bfseries R@3 & \bfseries R@5 & \bfseries MRR & \bfseries 失败模式 \\
\midrule
baseline hybrid        & 0.907 & 0.981 & 0.981 & 0.981 & — \\
\rowcolor{Cgraylt}
+reranker(chunk)       & 0.907 & 0.981 & 0.981 & \bad{0.975} & 子 chunk 丢上下文 \\
+reranker(parent, no dedup) & 0.907 & \bad{0.907} & \bad{0.907} & \bad{0.963} & \bad{同商品重复过多，挤掉其他候选} \\
\rowcolor{Cgraylt}
+reranker(parent+dedup) & 0.907 & 0.981 & 0.981 & 0.981 & 追平基线，无提升 \\
\bottomrule
\end{tabularx}}

\vspace{0.22cm}
\begin{tcolorbox}[tcred]
\small 一阶 hybrid recall@5 已达 \textbf{0.981}，无 headroom。reranker 两类失败：
\cnum{1} 子 chunk 丢上下文（MRR 0.981→0.975）；
\cnum{2} 父 chunk 无去重时同商品兄弟刷屏（R@3 \textbf{0.981→0.907}，多样性崩塌）。
\textbf{评测驱动决策：不上线 reranker。}
\end{tcolorbox}

\takeaway{负结果也是结果。更重要的是知道瓶颈在哪里，以及为什么暂时不做。}
\end{frame}


%% ── P9 FailureMemory ────────────────────────────────────
\begin{frame}
\frametitle{FailureMemory — 定位真正瓶颈}
\framesubtitle{不靠直觉，靠历史失败案例}

\begin{tcolorbox}[tcblue]
\footnotesize 累积 SupportCase → \mono{FailureMemory.analyze()} 批量归纳失败模式，\textbf{直接产出修复优先级}：
\end{tcolorbox}

\vspace{0.12cm}
{\footnotesize\renewcommand{\arraystretch}{1.12}
\begin{tabularx}{\textwidth}{lXX}
\toprule
\bfseries Pattern & \bfseries 根因诊断 & \bfseries 产出修复目标 \\
\midrule
\rowcolor{Cgraylt}
\mono{price\_constraint\_ignored}  & dense embedding 无法捕捉数值约束 & \hi{Q3 → price\_filter} \\
\mono{compound\_query\_recall\_gap} & ``A和B'' 遗漏一半实体 & \hi{Q28 → compound\_decomp} \\
\rowcolor{Cgraylt}
\mono{stale\_data\_caution}        & KB 文档 \mono{updated\_at} 超期 & \hi{freshness guardrail} \\
\bottomrule
\multicolumn{3}{l}{\scriptsize 其余 pattern 与 KB gap 放在备注页：\mono{zero\_retrieval\_handoff}、\mono{handoff\_cluster}、\mono{retrieval\_blind\_spot}、\mono{stale\_doc}}
\end{tabularx}}

\vspace{0.12cm}
\begin{columns}[T]
\column{0.49\textwidth}
\begin{tcolorbox}[tcgray, title={\color{Ctextdark}FM 与 CM 的分工}]
\footnotesize\begin{itemize}\setlength\itemsep{0.12em}
  \item \textbf{FailureMemory}：离线批量分析，\\pattern detection，输出修复方向
  \item \textbf{CaseMemory}：在线 advisory，\\query-time 语义检索，输出先验 trace
  \item {\color{Cmuted}一个诊断，一个预防}
\end{itemize}
\end{tcolorbox}

\column{0.48\textwidth}
\begin{tcolorbox}[tcblue]
\footnotesize FM 的输出是「该修哪里」；\\CM 的输出是「这题历史上有没有类似失败，给个 hint」。\\[0.15em]
两者互补，共同构成飞轮的学习闭环。
\end{tcolorbox}
\end{columns}

\takeaway{FailureMemory 的作用，是把失败整理成可执行的修复列表。}
\end{frame}


%% ── P10 两个精准修复 ─────────────────────────────────────
\begin{frame}
\frametitle{两个精准修复：Q3 + Q28}
\framesubtitle{针对 embedding 的已测失效模式}

\begin{columns}[T]
\column{0.49\textwidth}
\begin{tcolorbox}[tcblue, title={\color{Cblue}Q3 — 价格约束过滤}]
\small\begin{itemize}\setlength\itemsep{0.35em}
  \item \textbf{根因}：embedding ``预算600以内'' 仍把
    P009@\bad{899元}排名第一\\
    {\color{Cmuted}（数值约束 $\neq$ 语义相似度问题）}
  \item \textbf{修复}：\mono{price\_filter.parse\_budget()}\\
    解析中文预算 → post-retrieval 过滤 price $>$ budget
  \item \textbf{结果}：Q3 recall@1：\bad{0} → \hi{\textbf{1.0}}
\end{itemize}
\end{tcolorbox}

\column{0.49\textwidth}
\begin{tcolorbox}[tcgreen, title={\color{Cgreen}Q28 — 复合查询分解}]
\small\begin{itemize}\setlength\itemsep{0.35em}
  \item \textbf{根因}：``保温杯''子查询返回 P014\\
    {\color{Cmuted}（P014 Q\&A 含``和普通保温杯有什么不同''，语义污染）}
  \item \textbf{修复}：detect(``A和B'') → 分别检索 →\\
    doc 级 RRF → \textbf{标题 bigram 保底注入}（绕过语义歧义）
  \item \textbf{结果}：Q28 recall@5：\bad{0.5} → \hi{\textbf{1.0}}
\end{itemize}
\end{tcolorbox}
\end{columns}

\vspace{0.22cm}
\begin{tcolorbox}[tcgray]
\small\textbf{核心定位}：这些修复都来自 gold-doc 评测和 FailureMemory 的定位，属于\textbf{确定性修复层}：专门处理已经测出来的失效模式，影响范围清楚，也便于回滚。
\end{tcolorbox}

\takeaway{规则层的作用，是补足模型在特定场景下已经测出来的盲区。}
\end{frame}


%% ── P11 Ablation Table ──────────────────────────────────
\begin{frame}
\frametitle{Ablation Table}
\framesubtitle{最终量化证明（28 eval questions）}

{\footnotesize\renewcommand{\arraystretch}{1.18}
\begin{tabularx}{\textwidth}{lccccX}
\toprule
\bfseries 配置 & \bfseries R@1 & \bfseries R@3 & \bfseries R@5 & \bfseries MRR & \bfseries 说明 \\
\midrule
baseline               & 0.907 & 0.981 & 0.981 & 0.981 & 纯 hybrid \\
\rowcolor{Cgraylt}
+price\_filter         & \hi{0.944} & 0.981 & 0.981 & \hi{1.000} & $+\;0.037$ on R@1,\ $+\;0.019$ on MRR \\
+compound\_decomp      & 0.907 & 0.981 & \hi{1.000} & 0.981 & $\Delta$r@5=+0.019 \\
\rowcolor{Cgreelt}
\bfseries +pf+cd（final retrieval） & \hi{\bfseries 0.944} & \bfseries 0.981 & \hi{\bfseries 1.000} & \hi{\bfseries 1.000} & 两项可叠加，\hi{无回退} \\
+freshness\_guard      & 0.944 & 0.981 & 1.000 & 1.000 & 检索不变；guardrail 在生成层 \\
\rowcolor{Cgraylt}
memory\_prior          & — & — & — & — & advisory trace，非检索项 \\
\bottomrule
\end{tabularx}}

\vspace{0.12cm}
\begin{tcolorbox}[tcblue]
\footnotesize
price filter 修复的是\textbf{排名错误}（P009 错误占据第一）；  
compound decomposition 修复的是\textbf{候选覆盖缺口}（P006 根本不在 top-5）。  
两类不同的召回失败，对应两种定向修复，结果可以叠加，而且没有引入回退。
\end{tcolorbox}

\takeaway{每个模块贡献独立可测，两项核心修复可加且无回退——这是工程可信度的来源。}
\end{frame}


%% ── P12 记忆飞轮 ────────────────────────────────────────
\section{系统成熟度}

\begin{frame}
\frametitle{SupportCase + 记忆飞轮}
\framesubtitle{失败驱动进化}

\begin{columns}[T]
\column{0.50\textwidth}
\begin{tcolorbox}[tcgray, title={\color{Ctextdark}飞轮流程}]
{\fontsize{7.5}{10}\selectfont\ttfamily\color{Ctextmain}
handoff / 低置信案例\\[0.15em]
\quad$\downarrow$ SupportCase 落库 (SQLite, id=sc\_\{ts\}\_\{hash10\})\\[0.1em]
\quad$\downarrow$ FailureMemory.analyze()\\
\quad\quad $\rightarrow$ pattern + KB 缺口\\[0.1em]
\quad$\downarrow$ CaseMemory.build()\\
\quad\quad $\rightarrow$ embed 23 cases $\rightarrow$ .npy (threshold 0.70)\\[0.1em]
\quad$\downarrow$ 新 query 命中 (sim = 0.701) $\rightarrow$ trace\\[0.1em]
\quad$\downarrow$ suggested\_action = apply\_price\_filter\\
\quad\quad $\rightarrow$ 工程修复优先级\\
\quad\quad $\rightarrow$ price\_filter 上线 + ablation 验证
}
\end{tcolorbox}

\column{0.46\textwidth}
\begin{tcolorbox}[tcpurple, title={\color{purple!60}实测 trace 示例}]
{\fontsize{8}{10}\selectfont\ttfamily\color{purple!80!black}
记忆先验：sim=0.701\\
pattern=price\_constraint\_ignored\\
suggested=apply\_price\_filter
}
\end{tcolorbox}
\vspace{0.2cm}
\begin{tcolorbox}[tcblue]
\small
\textbf{重要说明}：CaseMemory 只提供 advisory trace，不自动执行修复；真正的修复还是靠工程实现和 ablation 验证完成。\\[0.2em]
更接近 episodic memory（语义检索）而非 cache（精确命中），threshold 0.70 控制误匹配。
\end{tcolorbox}
\end{columns}

\takeaway{系统不只记录回答，也记录失败和后续修复线索。}
\end{frame}


%% ── P13 新鲜度护栏 ──────────────────────────────────────
\begin{frame}
\frametitle{新鲜度护栏}
\framesubtitle{可信比「答得像」更重要}

\begin{columns}[T]
\column{0.50\textwidth}
\begin{tcolorbox}[tcgray, title={\color{Ctextdark}机制}]
\small\begin{itemize}\setlength\itemsep{0.3em}
  \item answer 中的价格/库存/政策 → 追溯\\chunk 的 \mono{updated\_at}
  \item 与 \mono{FRESHNESS\_MAX\_AGE\_DAYS=30} 对比
  \item 状态：\hi{fresh} / {\bfseries\color{orange}unverified} / \bad{stale}
  \item stale/unverified → 降级 ok → {\bfseries\color{orange}caution}，追加 advisory note
\end{itemize}
\end{tcolorbox}
\vspace{0.2cm}
\begin{tcolorbox}[tcgray, title={\color{Ctextdark}演示}]
\small\begin{itemize}\setlength\itemsep{0.2em}
  \item P005（\mono{updated\_at=2026-01-01}，故意留旧）→ \bad{stale caution}
  \item P006（\mono{2026-06-10}）→ \hi{fresh，正常通过}
\end{itemize}
\end{tcolorbox}

\column{0.46\textwidth}
\begin{tcolorbox}[tcgreen, title={\color{Cgreen}Ablation 验证}]
\small
+freshness\_guard 检索指标不变\\
（guardrail 在生成层，不影响召回）
\end{tcolorbox}
\vspace{0.2cm}
\begin{tcolorbox}[tcblue]
\small\textbf{关键说明}：freshness 的成功标准，不是 recall 提升，而是\textbf{避免过期价格或库存被当作确定事实输出}。\\[0.3em]
所以即使检索指标不变，这个 feature 仍然有价值，因为它管的是输出风险，不是召回能力。
\end{tcolorbox}
\end{columns}

\takeaway{客服系统不能乱答价格和库存——可信比「答得像」更重要。}
\end{frame}


%% ── P14 总结 + 边界声明 ─────────────────────────────────
\begin{frame}
\frametitle{总结 + 边界声明}

\begin{columns}[T]
\column{0.50\textwidth}
{\small\bfseries\color{Cmuted}证明了三件事}\\[0.1cm]

\begin{tcolorbox}[tcblue, title={\cnum{1} 会做检索工程}]
\small
hybrid(dense+BM25+RRF) + 父子文档 + 复合分解\\
recall@1：\textbf{0.907 → 0.944}
\end{tcolorbox}
\vspace{0.1cm}
\begin{tcolorbox}[tcgreen, title={\cnum{2} 会做科学评测}]
\small
gold 标注 + recall@k/MRR + ablation table\\
每一步改动有数字支撑
\end{tcolorbox}
\vspace{0.1cm}
\begin{tcolorbox}[tcamber, title={\cnum{3} 会做「不加某模块」的判断}]
\small
reranker A/B → 负结果 → 定位真瓶颈 → 精准修复
\end{tcolorbox}

\column{0.46\textwidth}
{\small\bfseries\color{Cmuted}边界声明（诚实是加分项）}\\[0.1cm]

\begin{tcolorbox}[tcred]
\footnotesize\begin{itemize}\setlength\itemsep{0.12em}
  \item 本地 demo + NSCC 离线评测，非生产系统
  \item 无 LLM key 时生成层降级为 catalog 模板
  \item 语料 40 商品，测试集 28 题，corpus 规模受限
  \item reranker 实测无提升，\textbf{未上线}
  \item CaseMemory 是 advisory trace，不自动执行修复
\end{itemize}
\end{tcolorbox}

\vspace{0.2cm}
\begin{tcolorbox}[tcblue]
\footnotesize
项目最终交付的，不只是一个 demo，\\
而是一套\textbf{可验证、可审计、支持持续迭代}的客服决策流程。
\end{tcolorbox}
\end{columns}

\takeaway{科学方法论驱动的 RAG 工程：量化评测 → 诚实负结果 → 定位真问题 → 修 → ablation 证明。}
\end{frame}


\end{document}
